\documentclass[11pt]{book}
\usepackage[paperwidth=145mm,paperheight=200mm,top=18mm,bottom=20mm,inner=20mm,outer=16mm]{geometry}
\usepackage{brumfiel}
% figures08.tex -- Chapter 8 (Parallel Lines) figures redrawn in TikZ.
%
% Local label clearance: ch08 has several figures where a vertex label sits
% at a point that two or three drawn lines pass through, so the shared
% 3.4pt outer sep is not enough.  Raised here only (brumfiel.sty untouched).
\tikzset{
  vlab/.style = {font=\footnotesize\itshape, inner sep=1.2pt, outer sep=5.2pt},
  slab/.style = {font=\scriptsize, inner sep=1.2pt, outer sep=5.2pt},
}
% Geometry read off the photo pages (book pp.133-153 = PDF 147-167) and the
% iPad scans (scan09 pp.7-23 = book pp.133-148; scan10 pp.1-4 = book
% pp.151-153).  Vertex ratios were measured off 400 dpi crops of the source
% figures; constrained points are constructed, never eyeballed: parallels via
% equal ratios, feet via projection, crossings via (intersection of ...).

% ---------------------------------------------------------------- Fig 8-1
% Two rays with the same direction: r1 || r2, both on the same side of the
% line through their vertices (Definition 8-2).
\newcommand{\FIGVIIIONE}{\begin{bkfigure}[0.62\linewidth]
\begin{tikzpicture}[scale=0.9]
  \coordinate (V1) at (0.62,1.15);          % vertex of r1
  \coordinate (V2) at (0.16,0.0);           % vertex of r2
  \draw[aux] ($(V2)!-0.30!(V1)$) -- ($(V2)!1.55!(V1)$);   % line of vertices
  \draw[fig,-latex] (V1) -- ++(4.35,0);
  \draw[fig,-latex] (V2) -- ++(4.35,0);
  \node[slab,above] at ($(V1)+(2.5,0)$) {$r_1$};
  \node[slab,above] at ($(V2)+(2.5,0)$) {$r_2$};
\end{tikzpicture}
\figcap{8--1}
\end{bkfigure}}

% ----------------------------------------------------------- Fig 8-2, 8-3
% 8-2: first proof of Thm 8-1.  n = AB; C on n beyond A; D on m and E on l
% chosen so that angle DAC is congruent to angle EBA, hence m || l, and the
% lines cannot meet at P.  Both congruent angles are arced, as in the book.
% 8-3: the lemma's figure -- the same three lines with m NOT parallel to l,
% and the two corresponding angles numbered 1 and 2.
\newcommand{\FIGVIIITWOTHREE}{\begin{bkfigure}[\linewidth]
\begin{minipage}{0.48\linewidth}\centering
\begin{tikzpicture}[scale=0.88]
  \coordinate (B) at (1.1470,0.0000);      \coordinate (A) at (1.9220,1.4756);
  \coordinate (C) at (2.1993,2.0035);      % on n, beyond A
  \coordinate (D) at (3.0194,1.4756);      % on m
  \coordinate (E) at (2.5110,0.0000);      % on l
  \coordinate (P) at (5.3320,0.7626);
  \draw[fig] (0.4650,0) -- (3.8440,0);                        % l (solid part)
  \draw[fig] (0.9378,-0.3984) -- (2.3295,2.2514);             % n
  \draw[aux] (0.4340,1.4756) -- (3.3790,1.4756);              % m (dashed)
  \draw[aux] (3.3790,1.4756) .. controls +(0.62,-0.06) and +(-0.50,0.38) .. (P);
  \draw[aux] (3.8440,0)      .. controls +(0.72,0.06) and +(-0.66,-0.34) .. (P);
  \draw[fig] ($(P)+(0.16,0.20)$) -- (P) -- ($(P)+(0.16,-0.20)$);
  % the two congruent angles of the construction
  \draw[tick] ($(A)+(0:0.30)$)  arc (0:62.29:0.30);
  \draw[tick] ($(B)+(0:0.36)$)  arc (0:62.29:0.36);
  \foreach \p in {C,D,E}{\fill (\p) circle (1.0pt);}
  \node[vlab] at (1.5100,1.8700) {$A$};
  \node[vlab,right] at (C) {$C$};
  \node[vlab,above] at (D) {$D$};
  \node[vlab,above] at (E) {$E$};   % book sets E above its dot, as with D
  \node[vlab] at (1.5200,-0.3300) {$B$};
  \node[vlab,right] at (P) {$P$};
  \node[vlab,below right] at (D) {$m$};
  \node[vlab,above] at (1.75,0) {$l$};
  \node[vlab] at (1.1200,0.8500) {$n$};
\end{tikzpicture}
\figcap{8--2}
\end{minipage}\hfill
\begin{minipage}{0.48\linewidth}\centering
\begin{tikzpicture}[scale=0.88]
  \coordinate (A) at (1.5190,1.4756);   \coordinate (B) at (0.7750,0.0000);
  \draw[fig] (-0.0250,-0.0594) -- (3.1950,0.1796);            % l through B
  \draw[fig] (0.7990,1.6310)  -- (3.5690,1.0332);             % m through A
  \draw[fig] (0.5769,-0.3929) -- (1.7531,1.9399);             % n through A,B
  \draw[fig,-latex] (2.387,0.620) -- (2.945,0.620);
  \node[vlab,right] at (2.945,0.620) {$P$};
  \draw[tick] ($(A)+(-12.18:0.34)$) arc (-12.18:63.24:0.34);
  \draw[tick] ($(B)+(4.24:0.38)$)   arc (4.24:63.24:0.38);
  \node[slab] at (2.0604,1.7342) {$1$};
  \node[slab] at (1.3238,0.3666) {$2$};
  \node[vlab] at (0.9050,1.1830) {$A$};
  \node[vlab] at (0.1400,-0.4000) {$B$};
  \node[vlab,right] at (3.5690,1.0332) {$m$};
  \node[vlab,above] at (2.1000,0.0983) {$l$};
  \node[vlab] at (0.5000,0.7800) {$n$};
\end{tikzpicture}
\figcap{8--3}
\end{minipage}
\end{bkfigure}}

% ---------------------------------------------------------------- Fig 8-4
% Third proof: n cuts m at A and l at B, with the indicated angles at A and B
% congruent (both arced).  Q is on m with AQ congruent to BP.
% The book draws m and l PARALLEL (both horizontal) and bends each with a
% dashed curve to the hypothetical meeting point P -- the same convention as
% Fig 8-2.  That is what makes the two arced angles at A and B congruent, as
% the third proof requires; m must not be tilted off l.
\newcommand{\FIGVIIIFOUR}{\begin{bkfigure}[0.66\linewidth]
\begin{tikzpicture}[scale=0.86]
  \coordinate (B) at (2.0460,0.0000);  \coordinate (A) at (2.8830,1.6244);
  \coordinate (P) at (4.5880,0.9424);
  % Q on m, with m || l through A and AQ = BP = 2.7111 exactly
  % (both re-checked in tools/constraints/ch08.py)
  \coordinate (Q) at (0.1719,1.6244);
  \draw[fig] (0.0000,0) -- (B);                          % l, drawn up to B
  \draw[fig] (1.9269,-0.2311) -- (3.0342,1.9177);        % n
  \draw[aux] (Q) -- (A);                                 % m
  \draw[fig] (Q) -- (B);
  \draw[aux] (A) .. controls +(1.00,-0.07) and +(-0.52,0.40) .. (P);
  \draw[aux] (B) .. controls +(1.08,0.10) and +(-0.66,-0.42) .. (P);
  \draw[fig] ($(P)+(0.17,0.20)$) -- (P) -- ($(P)+(0.17,-0.20)$);
  % the two congruent angles at A and B (both 62.74 deg, since m || l)
  \draw[tick] ($(A)+(0:0.32)$) arc (0:62.74:0.32);
  \draw[tick] ($(B)+(0:0.36)$) arc (0:62.74:0.36);
  \node[vlab,above left] at (A) {$A$};
  \node[vlab,below right] at (B) {$B$};
  \node[vlab,above] at (Q) {$Q$};
  \node[vlab,right] at (P) {$P$};
  \node[vlab,left] at (Q) {$m$};
  \node[vlab,left] at (0.0000,0) {$l$};
  \node[vlab,above right] at (3.0342,1.9177) {$n$};
\end{tikzpicture}
\figcap{8--4}
\end{bkfigure}}

% ----------------------------------------------------------- Fig 8-5, 8-6
% 8-5: l || m and m || n -- all three drawn parallel -- with l and n bent
% (dashed) to meet at P, the contradiction.
% 8-6: l || m and transversal n through P on l.
\newcommand{\FIGVIIIFIVESIX}{\begin{bkfigure}[\linewidth]
\begin{minipage}{0.48\linewidth}\centering
\begin{tikzpicture}[scale=0.80]
  \draw[fig] (0.30,1.90) -- (2.55,1.90);                  % l
  \draw[fig] (0.30,1.12) -- (2.55,1.12);                  % n
  \draw[fig] (0.30,0.34) -- (2.55,0.34);                  % m
  \coordinate (P) at (3.60,1.50);
  \draw[aux] (2.55,1.90) -- (P);
  \draw[aux] (2.55,1.12) -- (P);
  \draw[fig] ($(P)+(0.17,0.20)$) -- (P) -- ($(P)+(0.17,-0.20)$);
  \node[vlab,left] at (0.30,1.90) {$l$};
  \node[vlab,left] at (0.30,1.12) {$n$};
  \node[vlab,left] at (0.30,0.34) {$m$};
  \node[vlab,right] at (P) {$P$};
\end{tikzpicture}
\figcap{8--5}
\end{minipage}\hfill
\begin{minipage}{0.48\linewidth}\centering
\begin{tikzpicture}[scale=0.80]
  \draw[fig] (0.20,1.70) -- (3.70,1.70);                  % l
  \draw[fig] (0.20,0.42) -- (3.70,0.42);                  % m
  \coordinate (P) at (1.72,1.70);
  \draw[fig] ($(P)!-0.28!(2.86,0.10)$) -- (2.86,0.10);    % n
  \fill (P) circle (1.0pt);
  \node[vlab,above right] at (P) {$P$};
  \node[vlab,right] at (3.70,1.70) {$l$};
  \node[vlab,right] at (3.70,0.42) {$m$};
  \node[vlab,below] at (2.86,0.10) {$n$};
\end{tikzpicture}
\figcap{8--6}
\end{minipage}
\end{bkfigure}}

% ----------------------------------------------------------- Fig 8-7, 8-8
% 8-7: angle ABC congruent to angle BCD, hence AB || CD.
% 8-8: AB = CD, angles at B and C right, hence AD || BC.
\newcommand{\FIGVIIISEVENEIGHT}{\begin{bkfigure}[\linewidth]
\begin{minipage}{0.50\linewidth}\centering
\begin{tikzpicture}[scale=0.80]
  \coordinate (A) at (0,0);    \coordinate (B) at (1.72,0);
  \coordinate (C) at (2.52,0.86);
  \coordinate (D) at ($(C)+(1.72,0)$);
  \draw[fig] (A) -- (B) -- (C) -- (D);
  \node[vlab,left] at (A) {$A$};
  \node[vlab,below right] at (B) {$B$};
  \node[vlab,above left] at (C) {$C$};
  \node[vlab,right] at (D) {$D$};
\end{tikzpicture}
\figcap{8--7}
\end{minipage}\hfill
\begin{minipage}{0.46\linewidth}\centering
\begin{tikzpicture}[scale=0.80]
  \coordinate (B) at (0,0);    \coordinate (C) at (3.05,0);
  \coordinate (A) at (0,1.72); \coordinate (D) at (3.05,1.72);
  \draw[fig] (A) -- (B) -- (C) -- (D) -- cycle;
  % right-angle boxes at B and C
  \draw[fig] ($(B)+(0,0.20)$) -- ($(B)+(0.20,0.20)$) -- ($(B)+(0.20,0)$);
  \draw[fig] ($(C)+(0,0.20)$) -- ($(C)+(-0.20,0.20)$) -- ($(C)+(-0.20,0)$);
  % double ticks on AB and DC
  \foreach \dy in {-0.07,0.07}{
    \draw[tick] ($(A)!0.5!(B)+(-0.13,\dy)$) -- ($(A)!0.5!(B)+(0.13,\dy)$);
    \draw[tick] ($(D)!0.5!(C)+(-0.13,\dy)$) -- ($(D)!0.5!(C)+(0.13,\dy)$);}
  \node[vlab,above left] at (A) {$A$};
  \node[vlab,below left] at (B) {$B$};
  \node[vlab,below right] at (C) {$C$};
  \node[vlab,above right] at (D) {$D$};
\end{tikzpicture}
\figcap{8--8}
\end{minipage}
\end{bkfigure}}

% --------------------------------------------------------- Fig 8-9, 8-10
% 8-9: transversal l cutting l1 and l2; the eight numbered angles.  As in the
% book, l leans the other way from l1/l2 -- its upper crossing is to the LEFT
% of its lower one.  Each number sits on the bisector of its own angle.
% 8-10: proof figure for Thm 8-4, with G chosen so that AG || l2 (drawn just
% off l1 on purpose -- the proof concludes that G must in fact lie on l1).
\newcommand{\FIGVIIININETEN}{\begin{bkfigure}[\linewidth]
\begin{minipage}{0.50\linewidth}\centering
\begin{tikzpicture}[scale=1.30]
  \coordinate (L1a) at (-0.0620,1.1036);  \coordinate (L1b) at (3.0380,1.0745);
  \coordinate (L2a) at ( 0.0620,0.1240);  \coordinate (L2b) at (3.0380,0.3887);
  \coordinate (La)  at ( 1.3640,1.5047);  \coordinate (Lb)  at (2.3355,0.0000);
  \draw[fig] (L1a) -- (L1b);                              % l1
  \draw[fig] (L2a) -- (L2b);                              % l2
  \draw[fig] (La) -- (Lb);                                % l
  \node[slab] at (1.8070,1.4030) {$1$};
  \node[slab] at (1.3180,1.2614) {$2$};
  \node[slab] at (1.4596,0.7724) {$3$};
  \node[slab] at (1.9486,0.9140) {$4$};
  \node[slab] at (2.2943,0.6320) {$5$};
  \node[slab] at (1.8128,0.4665) {$6$};
  \node[slab] at (1.9783,-0.0150) {$7$};
  \node[slab] at (2.4598,0.1505) {$8$};
  \node[vlab,right] at (L1b) {$l_1$};
  \node[vlab,right] at (L2b) {$l_2$};
  \node[vlab,above] at (La) {$l$};
\end{tikzpicture}
\figcap{8--9}
\end{minipage}\hfill
\begin{minipage}{0.46\linewidth}\centering
\begin{tikzpicture}[scale=0.92]
  \coordinate (A) at (1.62,1.42);  \coordinate (B) at (1.34,0.30);
  \draw[fig] (-0.20,1.42) -- (3.60,1.42);                 % l1
  \draw[fig] (-0.20,0.30) -- (3.60,0.30);                 % l2
  \coordinate (E) at ($(B)!1.68!(A)$);
  \draw[fig] ($(B)!-0.42!(A)$) -- (E);                    % l
  \coordinate (C) at (0.42,1.42);
  \coordinate (F) at (3.00,1.42);
  \coordinate (D) at (2.85,0.30);
  % G on the ray from A that the proof compares with l1
  \coordinate (G) at (2.60,1.86);
  \draw[aux] (A) -- (G);
  \foreach \p in {C,F,D,G,E}{\fill (\p) circle (1.0pt);}  % book dots E too
  \node[vlab,above] at (C) {$C$};
  \node[vlab] at (1.374,1.735) {$A$};
  \node[vlab,above] at (F) {$F$};
  \node[vlab] at (1.586,-0.015) {$B$};
  \node[vlab,above] at (D) {$D$};
  \node[vlab,above left] at (E) {$E$};
  \node[vlab,above left] at (G) {$G$};
  \node[vlab,right] at (3.60,1.42) {$l_1$};
  \node[vlab,right] at (3.60,0.30) {$l_2$};
  \node[vlab,below left] at ($(B)!-0.42!(A)$) {$l$};
\end{tikzpicture}
\figcap{8--10}
\end{minipage}
\end{bkfigure}}

% --------------------------------------------------------------- Fig 8-11
% Bowtie: AO = CO and BO = OD, so AB || CD and BC || AD.
\newcommand{\FIGVIIIELEVEN}{\begin{bkfigure}[0.56\linewidth]
\begin{tikzpicture}[scale=0.86]
  \coordinate (O) at (2.05,0.92);
  \coordinate (A) at (1.15,0);      \coordinate (C) at ($(A)!2!(O)$);
  \coordinate (B) at (0,0.72);      \coordinate (D) at ($(B)!2!(O)$);
  \draw[fig] (A) -- (C);
  \draw[fig] (B) -- (D);
  \draw[fig] (B) -- (A);
  \draw[fig] (C) -- (D);
  \node[vlab,below] at (A) {$A$};
  \node[vlab,left]  at (B) {$B$};
  \node[vlab,above] at (C) {$C$};
  \node[vlab,right] at (D) {$D$};
  \node[vlab] at (1.86,1.36) {$O$};
\end{tikzpicture}
\figcap{8--11}
\end{bkfigure}}

% --------------------------------------------------------------- Fig 8-12
% Adjacent angles AOB and BOC: OB interior to angle AOC.
\newcommand{\FIGVIIITWELVE}{\begin{bkfigure}[0.48\linewidth]
\begin{tikzpicture}[scale=0.90]
  \coordinate (O) at (0,0);
  \coordinate (A) at (76:2.15);
  \coordinate (B) at (48:1.80);
  \coordinate (C) at (-9:2.45);
  \draw[fig] (O) -- (A);
  \draw[fig] (O) -- (B);
  \draw[fig] (O) -- (C);
  \node[vlab,above] at (A) {$A$};
  \node[vlab,above right] at (B) {$B$};
  \node[vlab,right] at (C) {$C$};
  \node[vlab,left] at (O) {$O$};
\end{tikzpicture}
\figcap{8--12}
\end{bkfigure}}

% --------------------------------------------------------------- Fig 8-13
% Thm 8-5: l through A parallel to BC; D and E on l, one on each side of AB.
\newcommand{\FIGVIIITHIRTEEN}{\begin{bkfigure}[0.58\linewidth]
\begin{tikzpicture}[scale=0.88]
  \coordinate (B) at (0,0);  \coordinate (C) at (3.35,0);
  \coordinate (A) at (1.675,1.90);
  \draw[fig] (B) -- (C) -- (A) -- cycle;
  % l is the horizontal through A: same direction as BC by construction
  \coordinate (D) at ($(A)+(-1.15,0)$);
  \coordinate (E) at ($(A)+(1.15,0)$);
  \draw[fig] ($(A)+(-1.72,0)$) -- ($(A)+(1.72,0)$);
  \fill (D) circle (1.0pt); \fill (E) circle (1.0pt);
  \node[vlab,above] at (D) {$D$};
  \node[vlab,above] at (A) {$A$};
  \node[vlab,above] at (E) {$E$};
  \node[vlab,left]  at (B) {$B$};
  \node[vlab,right] at (C) {$C$};
  \node[vlab,right] at ($(A)+(1.72,0)$) {$l$};
\end{tikzpicture}
\figcap{8--13}
\end{bkfigure}}

% --------------------------------------------------------------- Fig 8-14
% Two triangles with AB congruent to A'B', angle A congruent to angle A',
% angle C congruent to angle C'.  The book sets the primed triangle up and to
% the right, overlapping the first in x, and mirrored in sense.
\newcommand{\FIGVIIIFOURTEEN}{\begin{bkfigure}[0.66\linewidth]
\begin{tikzpicture}[scale=0.95]
  \coordinate (A) at (0,1.35);  \coordinate (B) at (0.39,0);
  \coordinate (C) at (2.34,0.78);
  \draw[fig] (A) -- (B) -- (C) -- cycle;
  \node[vlab,left] at (A) {$A$};
  \node[vlab,below] at (B) {$B$};
  \node[vlab,below right] at (C) {$C$};
  % A'B'C' solved from the three side lengths of ABC, so it is congruent to
  % ABC exactly (opposite sense, as the book draws it); checked in ch08.py.
  \coordinate (Cp) at (1.2000,2.1000);
  \coordinate (Ap) at (3.4950,2.8302);
  \coordinate (Bp) at (3.1992,1.4565);
  \draw[fig] (Ap) -- (Bp) -- (Cp) -- cycle;
  \node[vlab,above right] at (Ap) {$A'$};
  \node[vlab,below right] at (Bp) {$B'$};
  \node[vlab,left] at (Cp) {$C'$};
\end{tikzpicture}
\figcap{8--14}
\end{bkfigure}}

% -------------------------------------------------------- Fig 8-15, 8-16
% 8-15: angle A right, AD perpendicular to BC (D is the foot).
% 8-16: D the midpoint of BC, with BD congruent to AD.
\newcommand{\FIGVIIIFIFTEENSIXTEEN}{\begin{bkfigure}[\linewidth]
\begin{minipage}{0.48\linewidth}\centering
\begin{tikzpicture}[scale=0.80]
  % right angle at A: put A on the circle with diameter BC
  \coordinate (B) at (0,0);  \coordinate (C) at (3.45,0);
  \coordinate (M) at ($(B)!0.5!(C)$);
  \coordinate (A) at ($(M)+(65:1.725)$);
  \coordinate (D) at ($(B)!(A)!(C)$);       % foot of the altitude
  \draw[fig] (B) -- (C) -- (A) -- cycle;
  \draw[fig] (A) -- (D);
  \node[vlab,above] at (A) {$A$};
  \node[vlab,left]  at (B) {$B$};
  \node[vlab,right] at (C) {$C$};
  \node[vlab,below] at (D) {$D$};
\end{tikzpicture}
\figcap{8--15}
\end{minipage}\hfill
\begin{minipage}{0.48\linewidth}\centering
\begin{tikzpicture}[scale=0.80]
  \coordinate (B) at (0,0);  \coordinate (C) at (3.45,0);
  \coordinate (D) at ($(B)!0.5!(C)$);       % midpoint of BC
  % A on the circle centred D through B, so that BD = AD
  \coordinate (A) at ($(D)+(62:1.725)$);
  \draw[fig] (B) -- (C) -- (A) -- cycle;
  \draw[fig] (A) -- (D);
  \node[vlab,above] at (A) {$A$};
  \node[vlab,left]  at (B) {$B$};
  \node[vlab,right] at (C) {$C$};
  \node[vlab,below] at (D) {$D$};
\end{tikzpicture}
\figcap{8--16}
\end{minipage}
\end{bkfigure}}

% -------------------------------------------------------- Fig 8-17, 8-18
% 8-17: a polygonal path P1 ... Pn.  8-18: a six-sided path that crosses itself.
\newcommand{\FIGVIIISEVENTEENEIGHTEEN}{\begin{bkfigure}[\linewidth]
\begin{minipage}{0.48\linewidth}\centering
\begin{tikzpicture}[scale=0.74]
  \coordinate (P1) at (0,0.30);
  \coordinate (P2) at (1.15,0.92);
  \coordinate (P3) at (0.62,2.62);
  \coordinate (Q1) at (1.72,1.72);
  \coordinate (Q2) at (2.35,2.42);
  \coordinate (Q3) at (2.95,1.05);
  \coordinate (Pn) at (2.28,0);
  \draw[fig] (P1) -- (P2) -- (P3) -- (Q1) -- (Q2) -- (Q3) -- (Pn);
  \node[vlab,left]  at (P1) {$P_1$};
  \node[vlab,below right] at (P2) {$P_2$};
  \node[vlab,above left] at (P3) {$P_3$};
  \node[vlab,below right] at (Pn) {$P_n$};
\end{tikzpicture}
\figcap{8--17}
\end{minipage}\hfill
\begin{minipage}{0.48\linewidth}\centering
\begin{tikzpicture}[scale=0.74]
  \coordinate (P1) at (0,0.62);
  \coordinate (P2) at (3.35,1.42);
  \coordinate (P3) at (1.05,2.62);
  \coordinate (P4) at (0.62,0);
  \coordinate (P5) at (1.95,2.62);
  \coordinate (P6) at (2.05,0);
  \coordinate (P7) at (0.30,1.95);
  \draw[fig] (P1) -- (P2) -- (P3) -- (P4) -- (P5) -- (P6) -- (P7);
  \node[vlab,left]  at (P1) {$P_1$};
  \node[vlab,right] at (P2) {$P_2$};
  \node[vlab,above] at (P3) {$P_3$};
  \node[vlab,below] at (P4) {$P_4$};
  \node[vlab,above] at (P5) {$P_5$};
  \node[vlab,below] at (P6) {$P_6$};
  \node[vlab,left]  at (P7) {$P_7$};
\end{tikzpicture}
\figcap{8--18}
\end{minipage}
\end{bkfigure}}

% --------------------------------------------------------------- Fig 8-19
% Three simple polygons: a quadrilateral, a triangle, and a non-convex
% pentagon (the notch at P2).
\newcommand{\FIGVIIININETEEN}{\begin{bkfigure}[\linewidth]
\begin{minipage}{0.31\linewidth}\centering
\begin{tikzpicture}[scale=0.66]
  % measured off the book crop: both uprights vertical, top edge rising to
  % the right (the earlier version had it falling to the right)
  \coordinate (P1) at (0,0);       \coordinate (P4) at (2.350,0.044);
  \coordinate (P3) at (2.350,1.862); \coordinate (P2) at (0,1.703);
  \draw[fig] (P1) -- (P4) -- (P3) -- (P2) -- cycle;
  \node[vlab,below left] at (P1) {$P_1$};
  \node[vlab,above left] at (P2) {$P_2$};
  \node[vlab,above right] at (P3) {$P_3$};
  \node[vlab,below right] at (P4) {$P_4$};
\end{tikzpicture}
\end{minipage}\hfill
\begin{minipage}{0.31\linewidth}\centering
\begin{tikzpicture}[scale=0.66]
  \coordinate (P1) at (0,0);   \coordinate (P3) at (2.62,0);
  \coordinate (P2) at (1.01,1.82);   % apex at 0.386 of the base, as measured
  \draw[fig] (P1) -- (P3) -- (P2) -- cycle;
  \node[vlab,below left] at (P1) {$P_1$};
  \node[vlab,above] at (P2) {$P_2$};
  \node[vlab,below right] at (P3) {$P_3$};
\end{tikzpicture}
\end{minipage}\hfill
\begin{minipage}{0.31\linewidth}\centering
\begin{tikzpicture}[scale=0.66]
  % measured off the book crop: P1P2 nearly level, notch (reflex vertex) at P2
  \coordinate (P1) at (0,1.275);   \coordinate (P2) at (1.214,1.213);
  \coordinate (P3) at (1.565,2.418); \coordinate (P4) at (2.621,1.803);
  \coordinate (P5) at (1.346,0);
  \draw[fig] (P1) -- (P2) -- (P3) -- (P4) -- (P5) -- cycle;
  \node[vlab,left] at (P1) {$P_1$};
  % book sets P2 below-right of the notch; placed explicitly so the glyph
  % clears side P4P5, which runs close by on that side
  \node[vlab] at (1.250,0.700) {$P_2$};
  \node[vlab,above] at (P3) {$P_3$};
  \node[vlab,right] at (P4) {$P_4$};
  \node[vlab,below] at (P5) {$P_5$};
\end{tikzpicture}
\end{minipage}
\figcap{8--19.\ \ Simple polygons.}
\end{bkfigure}}

% -------------------------------------------------------- Fig 8-20, 8-21
% 8-20: naming parts of a polygon -- a quadrilateral with two consecutive
% angles arced, and (as the book draws it) a PENTAGON carrying two of its
% diagonals, dashed, with a pair of adjacent sides called out.
% 8-21: convex and non-convex.
\newcommand{\FIGVIIITWENTYTWENTYONE}{\begin{bkfigure}[\linewidth]
\begin{minipage}{0.60\linewidth}\centering
\begin{tikzpicture}[scale=0.836]
  % left quadrilateral: two consecutive angles marked with arcs
  \coordinate (a1) at (0,0);       \coordinate (a2) at (2.42,0);
  \coordinate (a3) at (2.42,1.42); \coordinate (a4) at (0.55,1.42);
  \draw[fig] (a1) -- (a2) -- (a3) -- (a4) -- cycle;
  \draw[tick] ($(a3)+(180:0.36)$) arc (180:270:0.36);
  \draw[tick] ($(a2)+(180:0.36)$) arc (180:90:0.36);
  \node[slab,align=center] at (1.15,2.70) {Consecutive\\angles};
  \draw[fig,-latex,line width=0.4pt] (1.20,2.00) -- (2.14,1.19);
  \draw[fig,-latex,line width=0.4pt] (1.34,2.00) -- (2.14,0.29);
  % right pentagon: adjacent sides and two diagonals
  \begin{scope}[shift={(3.10,0)}]
    \coordinate (b1) at (0,0);           % bottom-left
    \coordinate (b2) at (0.373,1.779);   % top-left
    \coordinate (b3) at (2.315,1.779);   % top-right
    \coordinate (b4) at (2.929,0.886);   % right point
    \coordinate (b5) at (2.286,0);       % bottom-right
    \draw[fig] (b1) -- (b2) -- (b3) -- (b4) -- (b5) -- cycle;
    \draw[aux] (b1) -- (b3);
    \draw[aux] (b3) -- (b5);
    \node[slab,align=center] at (0.35,2.70) {Adjacent\\sides};
    \node[slab] at (2.45,2.86) {Diagonals};
    \draw[fig,-latex,line width=0.4pt] (0.24,2.00) -- (0.131,0.623);
    \draw[fig,-latex,line width=0.4pt] (0.62,2.00) -- (0.850,1.779);
    \draw[fig,-latex,line width=0.4pt] (2.18,2.40) -- (1.572,1.208);
    \draw[fig,-latex,line width=0.4pt] (2.62,2.40) -- (2.301,0.943);
  \end{scope}
\end{tikzpicture}
\figcap{8--20}
\end{minipage}\hfill
\begin{minipage}{0.36\linewidth}\centering
\begin{tikzpicture}[scale=0.62]
  \coordinate (c1) at (0,0);   \coordinate (c2) at (1.86,0);
  \coordinate (c3) at (2.05,1.42); \coordinate (c4) at (0.42,1.62);
  \draw[fig] (c1) -- (c2) -- (c3) -- (c4) -- cycle;
  \node[slab] at (1.05,2.20) {Convex};
  \begin{scope}[shift={(2.75,0)}]
    \coordinate (d1) at (0,0.55);  \coordinate (d2) at (1.30,1.05);
    \coordinate (d3) at (2.62,0);  \coordinate (d4) at (1.15,0.62);
    \draw[fig] (d1) -- (d2) -- (d3) -- (d4) -- cycle;
    \node[slab] at (1.30,2.20) {Not convex};
  \end{scope}
\end{tikzpicture}
\figcap{8--21}
\end{minipage}
\end{bkfigure}}

% -------------------------------------------------------- Fig 8-22, 8-23
% 8-22: a pentagon cut into three triangles by the two diagonals (dashed)
% from one vertex, with all five exterior angles extended (dashed) and arced.
% 8-23: the same idea for an n-gon -- the book draws a seven-gon split by the
% four diagonals from one vertex, with two exterior angles suggested.
\newcommand{\FIGVIIITWENTYTWOTHREE}{\begin{bkfigure}[\linewidth]
\begin{minipage}{0.48\linewidth}\centering
\begin{tikzpicture}[scale=0.88]
  % Vertex ratios measured off a 400 dpi crop of the book figure with a 50 px
  % grid (book p.147): side v1v2 horizontal, side v2v3 vertical, v4 nearly
  % above v1.  The earlier coordinates were noticeably taller and narrower.
  \coordinate (v1) at (1.3360,0.0000); \coordinate (v2) at (2.9500,0.0230);
  \coordinate (v3) at (2.9090,1.0060); \coordinate (v4) at (1.0290,1.6310);
  \coordinate (v5) at (0.0000,0.5490);
  \draw[fig] (v1) -- (v2) -- (v3) -- (v4) -- (v5) -- cycle;
  \draw[aux] (v1) -- (v3);
  \draw[aux] (v1) -- (v4);
  % exterior-angle extensions: at each vertex, extend the side arrived on
  \draw[aux] (v1) -- (1.7060,-0.1520);
  \draw[aux] (v2) -- (3.3500,0.0287);
  \draw[aux] (v3) -- (2.8923,1.4057);
  \draw[aux] (v4) -- (0.6494,1.7572);
  \draw[aux] (v5) -- (-0.2757,0.2591);
  \draw[tick] ($(v1)+(337.66:0.26)$) arc (337.66:360.82:0.26);
  \draw[tick] ($(v2)+(0.82:0.26)$)   arc (0.82:92.39:0.26);
  \draw[tick] ($(v3)+(92.39:0.26)$)  arc (92.39:161.61:0.26);
  \draw[tick] ($(v4)+(161.61:0.26)$) arc (161.61:226.44:0.26);
  \draw[tick] ($(v5)+(226.44:0.26)$) arc (226.44:337.66:0.26);
  % each number sits on its exterior angle's bisector, at a radius chosen so
  % the glyph clears both bounding rays (the angle at v1 is only 23 deg, so
  % that one has to sit further out)
  \node[slab] at (2.1711,-0.1587) {$1$};
  \node[slab] at (3.2935,0.3863)  {$2$};
  \node[slab] at (2.6081,1.4053)  {$3$};
  \node[slab] at (0.5245,1.5050)  {$4$};
  \node[slab] at (0.1044,0.0600)  {$5$};
\end{tikzpicture}
\figcap{8--22.\ \ Pentagon\\(three triangles).}
\end{minipage}\hfill
\begin{minipage}{0.48\linewidth}\centering
\begin{tikzpicture}[scale=0.88]
  \coordinate (w1) at (0.0000,1.6150);   % the vertex the diagonals come from
  \coordinate (w2) at (0.3610,2.6980);
  \coordinate (w3) at (0.8740,2.7930);
  \coordinate (w4) at (3.0400,1.4630);
  \coordinate (w5) at (2.8690,0.6840);
  \coordinate (w6) at (1.9190,0.0000);
  \coordinate (w7) at (0.7600,0.0380);
  \draw[fig] (w1) -- (w2) -- (w3) -- (w4) -- (w5) -- (w6) -- (w7) -- cycle;
  \foreach \v in {w3,w4,w5,w6}{\draw[fig] (w1) -- (\v);}
  \draw[aux] (w3) -- (1.4836,2.9059);
  \draw[aux] (w4) -- (3.1729,2.0686);
\end{tikzpicture}
\figcap{8--23.\ \ $n$-gon\\($n{-}2$ triangles).}
\end{minipage}
\end{bkfigure}}

% --------------------------------------------------------------- Fig 8-24
% Regular five-pointed star: the points are the vertices of a regular pentagon.
\newcommand{\FIGVIIITWENTYFOUR}{\begin{bkfigure}[0.42\linewidth]
\begin{tikzpicture}[scale=0.90]
  \foreach \i in {0,...,4}{\coordinate (s\i) at ({90+72*\i}:1.72);}
  \draw[fig] (s0) -- (s2) -- (s4) -- (s1) -- (s3) -- cycle;
\end{tikzpicture}
\figcap{8--24}
\end{bkfigure}}

% --------------------------------------------------------------- Fig 8-25
% Parallelogram, rectangle, square, rhombus.
\newcommand{\FIGVIIITWENTYFIVE}{\begin{bkfigure}[\linewidth]
\begin{tikzpicture}[scale=0.70]
  % parallelogram
  \begin{scope}
    \draw[fig] (0,0) -- (1.42,0) -- (1.86,1.72) -- (0.44,1.72) -- cycle;
    \node[slab,below] at (0.92,-0.12) {Parallelogram};
  \end{scope}
  % rectangle
  \begin{scope}[shift={(3.05,0)}]
    \draw[fig] (0,0) -- (2.05,0) -- (2.05,1.62) -- (0,1.62) -- cycle;
    \node[slab,below] at (1.02,-0.12) {Rectangle};
  \end{scope}
  % square
  \begin{scope}[shift={(6.30,0)}]
    \draw[fig] (0,0) -- (1.62,0) -- (1.62,1.62) -- (0,1.62) -- cycle;
    \node[slab,below] at (0.80,-0.12) {Square};
  \end{scope}
  % rhombus
  \begin{scope}[shift={(8.95,0)}]
    \draw[fig] (0,0) -- (1.62,0) -- (2.06,1.56) -- (0.44,1.56) -- cycle;
    \node[slab,below] at (1.02,-0.12) {Rhombus};
  \end{scope}
\end{tikzpicture}
\figcap{8--25}
\end{bkfigure}}

% -------------------------------------------------------- Fig 8-26, 8-27
% 8-26: DE || AC and DF || AB.   8-27: regular hexagon ABCDEF, diagonal FD.
\newcommand{\FIGVIIITWENTYSIXSEVEN}{\begin{bkfigure}[\linewidth]
\begin{minipage}{0.50\linewidth}\centering
\begin{tikzpicture}[scale=0.80]
  \coordinate (B) at (0,0);  \coordinate (C) at (3.45,0);
  \coordinate (A) at (1.72,1.86);
  \coordinate (D) at ($(B)!0.52!(C)$);
  % E on AB and F on AC with DE || AC, DF || AB  (equal ratios)
  \coordinate (E) at ($(B)!0.52!(A)$);
  \coordinate (F) at ($(C)!0.48!(A)$);
  \draw[fig] (B) -- (C) -- (A) -- cycle;
  \draw[aux] (D) -- (E);
  \draw[aux] (D) -- (F);
  \node[vlab,above] at (A) {$A$};
  \node[vlab,left]  at (B) {$B$};
  \node[vlab,right] at (C) {$C$};
  \node[vlab,below] at (D) {$D$};
  \node[vlab,above left]  at (E) {$E$};
  \node[vlab,above right] at (F) {$F$};
\end{tikzpicture}
\figcap{8--26}
\end{minipage}\hfill
\begin{minipage}{0.46\linewidth}\centering
\begin{tikzpicture}[scale=0.80]
  \foreach \i in {0,...,5}{\coordinate (h\i) at ({-60+60*\i}:1.42);}
  \draw[fig] (h0) -- (h1) -- (h2) -- (h3) -- (h4) -- (h5) -- cycle;
  \draw[aux] (h4) -- (h2);
  \node[vlab,below right] at (h0) {$B$};
  \node[vlab,right] at (h1) {$C$};
  \node[vlab,above right] at (h2) {$D$};
  \node[vlab,above left] at (h3) {$E$};
  \node[vlab,left] at (h4) {$F$};
  \node[vlab,below left] at (h5) {$A$};
\end{tikzpicture}
\figcap{8--27}
\end{minipage}
\end{bkfigure}}

% -------------------------------------------------------- Fig 8-28, 8-29
% 8-28: two cevians AQ and AR from A.  The lettered angles are, exactly as
% the book places them:  x = angle B, y = angle BAQ, s = angle RAC,
% z = angle AQC (right of AQ), t = angle ARB (left of AR), r = angle C.
% With that reading x + y + t = r + s + z, which is what Ex. 9 asks for; it
% is verified numerically in tools/constraints/ch08.py.
% 8-29: AB = BC, BE = BD, angle CBD = 40 degrees; angle B obtuse as drawn.
\newcommand{\FIGVIIITWENTYEIGHTNINE}{\begin{bkfigure}[\linewidth]
\begin{minipage}{0.48\linewidth}\centering
\begin{tikzpicture}[scale=1.05]
  % Ratios measured off a 400 dpi crop of the book figure (book p.151):
  % Q at 0.369 and R at 0.620 of BC, apex almost exactly over Q so that AQ
  % reads as vertical.  A is nudged 0.04 left of Q so that angle AQC is
  % 91.3 deg, as the book draws it, rather than an unstated exact right angle.
  \coordinate (B) at (0,0);  \coordinate (C) at (3.45,0);
  \coordinate (A) at (1.2331,1.7664);
  \coordinate (Q) at ($(B)!0.369!(C)$);
  \coordinate (R) at ($(B)!0.620!(C)$);
  \draw[fig] (B) -- (C) -- (A) -- cycle;
  \draw[fig] (A) -- (Q);
  \draw[fig] (A) -- (R);
  \node[vlab,above] at (A) {$A$};
  \node[vlab,left]  at (B) {$B$};
  \node[vlab,right] at (C) {$C$};
  % the two angles at A carry arcs, as in the book
  \draw[tick] ($(A)+(-124.92:0.45)$) arc (-124.92:-88.70:0.45);
  \draw[tick] ($(A)+(-62.85:0.50)$)  arc (-62.85:-38.55:0.50);
  % each letter on the bisector of the angle it names, far enough out to clear
  % both sides of that angle
  \node[slab] at (0.4079,0.2127) {$x$};
  \node[slab] at (1.0307,1.0963) {$y$};
  \node[slab] at (1.7716,1.1087) {$s$};
  \node[slab] at (1.5527,0.2860) {$z$};
  \node[slab] at (1.7465,0.2398) {$t$};
  \node[slab] at (2.8459,0.2113) {$r$};
\end{tikzpicture}
\figcap{8--28}
\end{minipage}\hfill
\begin{minipage}{0.48\linewidth}\centering
\begin{tikzpicture}[scale=0.80]
  \coordinate (B) at (0,0);  \coordinate (C) at (3.05,0);
  % AB = BC: A on the circle centred B with radius BC, at 110 deg as drawn
  \coordinate (A) at ($(B)+(110:3.05)$);
  % D is where the 40-degree ray from B meets AC.  Solved exactly and written
  % out as literals so no extra tikz library is needed; angle CBD = 40.0000
  % deg is re-checked in tools/constraints/ch08.py.
  \coordinate (D) at (1.38738,1.16417);
  % BE = BD: E on BA at the same distance from B as D  (ratio 0.59381)
  \coordinate (E) at ($(B)!0.59381!(A)$);
  \draw[fig] (B) -- (C) -- (A) -- cycle;
  \draw[fig] (B) -- (D);
  \draw[fig] (E) -- (D);
  \node[vlab,above left] at (A) {$A$};
  \node[vlab,below] at (B) {$B$};
  \node[vlab,below right] at (C) {$C$};
  \node[vlab] at (1.6670,1.4040) {$D$};
  \node[vlab] at (-1.3700,1.4300) {$E$};
  \node[slab] at (1.2592,0.4583) {$40^\circ$};
\end{tikzpicture}
\figcap{8--29}
\end{minipage}
\end{bkfigure}}

% --------------------------------------------------------------- Fig 8-30
% AB congruent to A'B', with B between A and C.
\newcommand{\FIGVIIITHIRTY}{\begin{bkfigure}[0.66\linewidth]
\begin{tikzpicture}[scale=0.90]
  \draw[fig] (-0.20,0) -- (2.20,0);
  \coordinate (A) at (0,0); \coordinate (B) at (1.05,0);
  \coordinate (C) at (1.86,0);
  \foreach \p in {A,B,C}{\fill (\p) circle (1.0pt);}
  \draw[tick] ($(A)!0.5!(B)+(0,-0.13)$) -- ($(A)!0.5!(B)+(0,0.13)$);
  \node[vlab,above] at (A) {$A$};
  \node[vlab,above] at (B) {$B$};
  \node[vlab,above] at (C) {$C$};
  \begin{scope}[shift={(3.05,0)}]
    \draw[fig] (-0.20,0) -- (1.60,0);
    \coordinate (Ap) at (0,0); \coordinate (Bp) at (1.05,0);
    \fill (Ap) circle (1.0pt); \fill (Bp) circle (1.0pt);
    \draw[tick] ($(Ap)!0.5!(Bp)+(0,-0.13)$) -- ($(Ap)!0.5!(Bp)+(0,0.13)$);
    \node[vlab,above] at (Ap) {$A'$};
    \node[vlab,above] at (Bp) {$B'$};
  \end{scope}
\end{tikzpicture}
\figcap{8--30}
\end{bkfigure}}

% -------------------------------------------------------- Fig 8-31, 8-32
% 8-31: angles CAB and ABD right, CA = DB, M and N midpoints of CD and AB.
% 8-32: X and Y on one side of AB; rays AX and BY meet at P.
\newcommand{\FIGVIIITHIRTYONETWO}{\begin{bkfigure}[\linewidth]
\begin{minipage}{0.50\linewidth}\centering
\begin{tikzpicture}[scale=0.80]
  \coordinate (A) at (0,0);      \coordinate (B) at (3.05,0);
  \coordinate (C) at (0,1.62);   \coordinate (D) at (3.05,1.62);
  \coordinate (M) at ($(C)!0.5!(D)$);
  \coordinate (N) at ($(A)!0.5!(B)$);
  \draw[fig] (A) -- (B) -- (D) -- (C) -- cycle;
  \draw[aux] (M) -- (N);
  % right-angle boxes at A and B
  \draw[fig] ($(A)+(0,0.20)$) -- ($(A)+(0.20,0.20)$) -- ($(A)+(0.20,0)$);
  \draw[fig] ($(B)+(0,0.20)$) -- ($(B)+(-0.20,0.20)$) -- ($(B)+(-0.20,0)$);
  % single ticks on CM, MD ; double ticks on AN, NB
  \draw[tick] ($(C)!0.5!(M)+(0,-0.13)$) -- ($(C)!0.5!(M)+(0,0.13)$);
  \draw[tick] ($(M)!0.5!(D)+(0,-0.13)$) -- ($(M)!0.5!(D)+(0,0.13)$);
  \foreach \dx in {-0.06,0.06}{
    \draw[tick] ($(A)!0.5!(N)+(\dx,-0.13)$) -- ($(A)!0.5!(N)+(\dx,0.13)$);
    \draw[tick] ($(N)!0.5!(B)+(\dx,-0.13)$) -- ($(N)!0.5!(B)+(\dx,0.13)$);}
  \node[vlab,above left] at (C) {$C$};
  \node[vlab,above] at (M) {$M$};
  \node[vlab,above right] at (D) {$D$};
  \node[vlab,below left] at (A) {$A$};
  \node[vlab,below] at (N) {$N$};
  \node[vlab,below right] at (B) {$B$};
\end{tikzpicture}
\figcap{8--31}
\end{minipage}\hfill
\begin{minipage}{0.46\linewidth}\centering
\begin{tikzpicture}[scale=0.80]
  \coordinate (A) at (0,1.72);  \coordinate (B) at (0,0);
  \coordinate (P) at (3.05,1.02);
  \coordinate (X) at ($(A)!0.55!(P)$);
  \coordinate (Y) at ($(B)!0.55!(P)$);
  \draw[fig] (A) -- (B);
  \draw[fig] (A) -- (X);
  \draw[fig] (B) -- (Y);
  \draw[aux] (X) -- (P);
  \draw[aux] (Y) -- (P);
  \draw[fig] ($(P)+(0.17,0.20)$) -- (P) -- ($(P)+(0.17,-0.20)$);
  \fill (X) circle (1.0pt); \fill (Y) circle (1.0pt);
  \node[vlab,above] at (A) {$A$};
  \node[vlab,below] at (B) {$B$};
  \node[vlab,above] at (X) {$X$};
  \node[vlab,anchor=north] at ($(Y)+(0,-0.32)$) {$Y$};
  \node[vlab,above right] at (P) {$P$};
\end{tikzpicture}
\figcap{8--32}
\end{minipage}
\end{bkfigure}}


% ---- local macros (hoist into book preamble) ----
% Halftone plate, never redrawn.  Carries the real image since 2026-08-17
% (Caleb: "I WANT THE ACTUAL HALFTONE PICTURES IN THERE"), cut from scan09
% pp.12-13 -- see plates/MANIFEST.md.
%   #1 = image file in plates/, #2 = caption block, #3 = figure number
\newcommand{\VIIIplate}[3]{%
  \begin{bkfigure}[0.78\linewidth]
  \includegraphics[height=62mm]{#1}\par\smallskip
  {\footnotesize #2}
  \figcap{#3}
  \end{bkfigure}}
% Two unnumbered "*" footnotes belonging to this chapter (book pp.135 and 139).
% \starnote in brumfiel.sty emits a bare \footnotetext, which prints the
% footnote counter; the book marks both of these with "*" and no number, so
% they get the same counter-suppressing treatment ch10 gives \Xlimitnote.
\newcommand{\VIIIproofnote}{\begingroup\renewcommand{\thefootnote}{}%
  \footnotetext{*\,If the teacher wishes, the second and third proofs may be
  omitted. Note that the second proof employs a lemma, or helping
  theorem.}\endgroup}
\newcommand{\VIIIgaussnote}{\begingroup\renewcommand{\thefootnote}{}%
  \footnotetext{*\,See, for example, the discussion of Gauss' experiment in
  \emph{The World of Mathematics}, by J.~R. Newman, Vol.~3, p.~1644.}\endgroup}
% The three-utilities puzzle on book p.145 is set as bare letters over dots in
% the running text, not as a numbered figure.
\newcommand{\VIIIutilities}{%
  \par\medskip\begin{center}
  \begin{tabular}{@{}c@{\hspace{3.2em}}c@{\hspace{3.2em}}c@{}}
    $A$ & $B$ & $C$ \\[-1pt]
    $\cdot$ & $\cdot$ & $\cdot$ \\[4pt]
    $G$ & $W$ & $E$ \\[-1pt]
    $\cdot$ & $\cdot$ & $\cdot$ \\
  \end{tabular}
  \end{center}\par\medskip}
% Narrow multicol columns plus this chapter's primed/subscripted math atoms
% ($\ang{A'O'B'}$, $l_1 \pll l_2$) leave TeX few legal break points.
\setlength{\emergencystretch}{2.5em}
% ---- end local macros ----

\begin{document}
\setcounter{chapter}{8}

\chapter*{\raggedright\Huge PARALLEL LINES}
\addcontentsline{toc}{chapter}{8.\ Parallel Lines}
\markboth{PARALLEL LINES}{}
\vspace{-1em}\noindent\rule{\linewidth}{0.6pt}\vspace{1em}
\figurekey

%========================================================= 8-1
\section*{8--1\quad DEFINITION OF PARALLEL LINES}
\addcontentsline{toc}{section}{8--1\quad Definition of parallel lines}

In earlier mathematics courses, you learned some facts about parallel lines.
Before reading this chapter, try to recall some properties of parallel lines
that you have accepted as true. The following exercises will help your memory
and prepare you for proofs.

\exercisegroup{8--1}
\begin{exercisecols}
\begin{exlist}
\item What do you mean when you say, ``Line $l$ is \emph{parallel} to line
  $m$''?
\item Name some geometric figures that have pairs of \emph{parallel} sides.
\item Explain why two sides of a triangle are not \emph{parallel}.
\item Two rocks are dropped, one from the Brooklyn Bridge, and the other from
  the Golden Gate Bridge. Would you say that these rocks fall along
  \emph{parallel} paths?
\item If you drop two rocks, one from each hand, do they traverse
  \emph{parallel} paths?
\item Roads $r$ and $s$ run ``due north and south.'' Are they \emph{parallel}?
\item Roads $t$ and $u$ run ``due east and west.'' Are they \emph{parallel}?
\item How many different pairs of \emph{parallel} edges does a cube have?
\item Both lines $l$ and $m \perp$ line $n$. Does this suggest a theorem about
  \emph{parallels}?
\item Each of the lines $l$ and $m$ is \emph{parallel} to line $n$. What
  theorem does this suggest?
\item Lines $p$ and $q$ are \emph{parallel}. Line $r \perp p$. What theorem
  does this suggest?
\item Draw two angles, $\ang{AOB}$ and $\ang{A'O'B'}$, such that the lines
  $AO$ and $A'O'$ are \emph{parallel}, and also lines $BO$ and $B'O'$ are
  \emph{parallel}. Does this figure suggest a theorem? Be careful! Draw more
  than one figure.
\item Draw a line $l$, and mark a point $A$ not on $l$. What might it mean to
  say that a ray from $A$ is \emph{parallel} to line $l$?
\item Which of the above exercises do not strictly belong to plane geometry?
\end{exlist}
\end{exercisecols}

If you discussed Exercises 6 and 7, someone probably pointed out that roads on
the surface of the earth are not quite the same as lines. Also, in connection
with Exercises 4 and 5, the paths along which objects fall are hardly the
imaginary lines of our geometry. Since the earth rotates about its axis and at
the same time revolves around the sun, an object falling toward the center of
the earth follows a path extremely hard to visualize. Our task is to choose a
definition of parallel lines, such that we can use the definition in our plane
geometry along with our postulates to prove the truth of statements that
concern parallel lines and that we feel ought to be true. The definition below
was perhaps suggested by someone in the class during the discussion of
Exercise Group 8--1.

\begin{definition}
Two lines $l$ and $m$ are parallel if and only if they do not intersect. We
indicate that two lines are parallel by writing ``$l \pll m$.''
\end{definition}

\exercisegroup{8--2}
\begin{exercisecols}
\begin{exlist}
\item Give definitions of the following. In later sections, it will be
  convenient to have these definitions. Good definitions are not hard to make,
  and so are given as exercises here.
  \begin{enumerate}[label=(\alph*), leftmargin=1.8em, topsep=2pt, itemsep=1pt]
    \item A ray parallel to a line.
    \item Parallel rays.
    \item A line segment parallel to a line.
    \item Parallel line segments.
  \end{enumerate}
\item Draw a figure of a line not intersecting a segment and not parallel to
  the segment.
\item Can a ray and a line not intersect and yet not be parallel?
\end{exlist}
\end{exercisecols}

Now that we have parallel lines, it is possible for us to talk about rays that
have the same direction. This requires another definition.

\begin{definition}
Two rays have the same \emph{direction} if and only if they are parallel and
both rays are on the same side of the line through their vertices
(Fig.~8--1).
\end{definition}

\FIGVIIIONE

\exercise{8--3}
\noindent What would have been wrong had we defined lines to be parallel if
they had the same direction?

%========================================================= 8-2
\section*{8--2\quad EXISTENCE OF PARALLELS}
\addcontentsline{toc}{section}{8--2\quad Existence of parallels}

Perhaps it will seem strange to you that we are going to prove that parallel
lines actually exist. We could add an extra postulate and assume that there
are parallels, but the proof is not difficult. We give three different proofs
of the following theorem.

\begin{theorem}
If $l$ is any line, and $A$ any point not on $l$, then there exists at least
one line through $A$ parallel to $l$.
\end{theorem}

\emph{First Proof.} Referring to Fig.~8--2, let $B$ be any point on $l$, and
let $n$ be the line through $A$ and $B$. Let $E$ be another point on $l$, and
let $D$ be a point on the same side of $n$ as $E$, such that
$\ang{DAC} \cong \ang{EBA}$. Lines $l$ and $m$ cannot intersect, for if they
did, say at $P$, we should have $\tri{PBA}$ with exterior
$\ang{PAC} \cong \ang{PBA}$. But we know this is impossible, for an exterior
angle of a triangle is always greater than either opposite interior angle.

\FIGVIIITWOTHREE

\emph{*Second Proof. Lemma.} Let line $l$, point $A$, and line $n$ be as in
Theorem 8--1.\VIIIproofnote{} Now, referring to Fig.~8--2, if $m$ is not
parallel to $l$, then $\ang 1$ is not congruent to $\ang 2$. (For ``not
parallel to,'' we use the symbol $\nparallel$.) \emph{Proof.} If $m$ is not
parallel to $l$, these lines intersect, a triangle is formed, and
$\ang 1 \ncong \ang 2$. Now Theorem 8--1 follows quickly from the lemma. The
lemma states that $l \nparallel m \imp \ang 1 \ncong \ang 2$. The
contrapositive of the lemma is $\ang 1 \cong \ang 2 \imp l \pll m$, and the
theorem obviously follows.

\emph{*Third Proof.} Referring to Fig.~8--4, let $l$, $m$, and $n$ refer to
the lines of the preceding proofs with $n$ cutting $m$ at $A$, and $l$ at $B$,
and the indicated angles at $A$ and $B$ congruent to each other. If $l$ and
$m$ intersected at $P$, we could choose $Q$ on $m$, as above, so that
$AQ \cong BP$. It follows at once, by the S.A.S. theorem, that
$\tri{BAQ} \cong \tri{ABP}$. Hence, $\ang{ABQ} \cong \ang{BAP}$.

\FIGVIIIFOUR

Since also, $\ang{ABP} \cong \ang{BAQ}$, we have
$\ang{QBP} \cong \ang{QAP}$. $\ang{QAP}$ is a straight angle, so $\ang{QBP}$
is a straight angle. It follows that $Q$ lies on line $l$, and so if lines $l$
and $m$ intersect in one point, they have a second point in common. This is
impossible, so lines $l$ and $m$ do not intersect.

\exercisegroup{8--4}
\begin{exercisecols}
\begin{exlist}
\item Draw a line, and mark a point not on it. Use the first proof of Theorem
  8--1 to draw a line, parallel to the first line, through your point.
\item Prove that if $l \perp m$, and $n \perp m$, then, if $l \neq n$,
  $l \pll n$.
\item Is a line parallel to itself?
\item In the first proof of Theorem 8--1, several details were omitted. First
  a point $B$ was chosen and then a point $E$. How was $C$ selected? How would
  we have argued if $l$ and $m$ were to intersect on the other side of $AB$?
\item The second proof of Theorem 8--1 is really just a reworded version of
  the first proof. Explain why this is so.
\item Without pointing at Fig.~8--4, describe how the point $Q$ is chosen in
  the third proof. What postulate guarantees the existence of $Q$?
\item Prove that if $l_1 \pll l_2$, then all the points of $l_2$ are on the
  same side of $l_1$.
\item Use the result of Exercise 7 to give a definition of ``the set of points
  between $l_1$ and $l_2$, where $l_1 \pll l_2$.''
\end{exlist}
\end{exercisecols}

%========================================================= 8-3
\section*{8--3\quad THE PARALLEL POSTULATE}
\addcontentsline{toc}{section}{8--3\quad The parallel postulate}

We proved in Section 8--2 that through a point $A$ not on line $l$, there is
\emph{at least one line parallel to $l$}. Now we raise the question, \emph{is
there more than one such line?} Our geometric intuition tells us that there is
only the one, and it is natural to attempt to prove this from our other
postulates. For 2000 years, many fine mathematicians tried and failed to find
such a proof. At last, in the 19th century, it was shown that \emph{it is
impossible to prove from the postulates we have already listed that there is
no more than one line through a given point parallel to a given line}. Euclid
himself assumed that there was no more than one such parallel, and we shall
make the same assumption.

\par\medskip\noindent
\textbf{Postulate VI--1.} \emph{The parallel postulate.} Through a point not
on a line $l$, there is no more than one line parallel to $l$.
\par\medskip

\emph{Historical note.} It may seem strange that it is possible to prove that
the parallel postulate is not a consequence of the other postulates. The
understanding of the fact that the parallel postulate is independent of the
others came about in the following way.

After vainly trying to prove the truth of the parallel postulate by using the
other postulates of Euclid, Johann Bolyai (Yo'-hahn Bull'-y\={\i}), a young
Hungarian mathematician, decided to assume that \emph{a second parallel could
be drawn through A}. Using this postulate, he proved many strange and
beautiful theorems. Later it was shown that Bolyai's geometry was just as
acceptable as ordinary Euclidean geometry. We call this geometry of Bolyai's a
\emph{non-Euclidean} geometry to contrast it with our Euclidean geometry which
employs the parallel postulate of Euclid. Whenever you encounter the term
non-Euclidean geometry, you can be sure that some postulate different from the
parallel postulate is used. All the theorems we have proved up to now are true
in the non-Euclidean geometry of Bolyai.

When these things are first explained, many students ask which postulate is
right. Can only one parallel be drawn, or is it possible to draw several?
While these questions are quite natural ones to ask, they do not have the
simple answers that one might expect. Our postulates are statements that we
invent in order to prove interesting and useful theorems. The theorems of
Euclidean geometry are useful in carpentry, drafting, surveying, navigation,
astronomy, and in many other fields. The theorems of some non-Euclidean
geometries have been found to be extremely useful in astronomy and physics.
Today mathematicians study many different geometries. Each of the geometries
has a postulate system different from that of any other geometry. All are
interesting. Many have proved to be of practical use. It does not make sense
to ask which geometry is true. Mathematicians design geometries as architects
design houses. You might as well ask which house is true. It \emph{is}
appropriate to ask which house is better adapted to one's living requirements.
And so it is also appropriate to ask which geometry is most convenient to
describe the universe.*\VIIIgaussnote

\VIIIplate{plate-gauss}{\textsc{Karl Friedrich Gauss}\\(1777--1855)\\
  Probably the first mathematician to recognize the\\
  existence of non-Euclidean geometries.\\
  Courtesy of \emph{Scripta Mathematica}.}{--- (book p.~137)}

\VIIIplate{plate-lobachevsky}{\textsc{Nicolai Ivanovich Lobachevsky}\\(1793--1856)\\
  Co-inventor, with Johann Bolyai, of\\
  non-Euclidean geometry.\\
  Courtesy of \emph{Scripta Mathematica}.}{--- (book p.~138)}

%========================================================= 8-4
\section*{8--4\quad CONSEQUENCES OF THE\\PARALLEL POSTULATE}
\addcontentsline{toc}{section}{8--4\quad Consequences of the parallel postulate}

We can now prove several interesting theorems quite easily.

\begin{theorem}
If $l$, $m$, and $n$ are distinct lines, with $l \pll m$ and $m \pll n$, then
$l \pll n$ \textup{(Fig.~8--5)}.
\end{theorem}

\emph{Proof.} If $l$ were not parallel to $n$, then $l$ and $n$ would
intersect, say at $P$ in Fig.~8--5. Two lines, namely $l$ and $n$, would then
pass through $P$, parallel to $m$, thus contradicting the parallel postulate.

\begin{theorem}
If $l \pll m$, and line $n \neq l$ intersects $l$ at $P$, then $n$ intersects
$m$ \textup{(Fig.~8--6)}.
\end{theorem}

\FIGVIIIFIVESIX

\noindent\emph{Proof.}
\begin{steps}
\step{$l \pll m$.}{Given.}
\step{$P$ is on $n$ and $l$, and $n \neq l$.}{Given.}
\step{$l$ is the only line on $P$ parallel to $m$.}{Statements 1, 2, and the
  parallel postulate.}
\step{$n$ intersects $m$.}{Statements 2 and 3.}
\end{steps}

\exercisegroup{8--5}
\begin{exercisecols}
\begin{exlist}
\item Prove that the construction of a parallel to $l$ through a given point
  $A$ is an \RC{} construction.
\item In Fig.~8--7, $\ang{ABC} \cong \ang{BCD}$. Prove that $AB \pll CD$.
\exfig{\FIGVIIISEVEN}
\item Prove that if $l \perp n$, $n \perp m$, and $l \neq m$, then
  $l \pll m$. Have you used the parallel postulate?
\item If $l \pll m$, and $l \perp n$, prove that $n \perp m$. Did you use the
  parallel postulate?
\item If $l \pll m$, $k \pll m$, and $l$ and $k$ intersect at $P$, what can
  you say?
\item[\stex 6.] In Fig.~8--8, $\sg{AB} = \sg{CD}$, and $\ang{ABC}$ and
  $\ang{BCD}$ are right. Prove that $AD \pll BC$. [\emph{Hint:} Prove that
  $\ang A \cong \ang D$.]
\exfig{\FIGVIIIEIGHT}
\item[\stex 7.] If a rectangle is defined as a quadrilateral with four right
  angles, prove that rectangles actually exist.
\end{exlist}
\end{exercisecols}

%========================================================= 8-5
\section*{8--5\quad CORRESPONDING AND\\ALTERNATE ANGLES}
\addcontentsline{toc}{section}{8--5\quad Corresponding and alternate angles}

When a line intersects two other lines, we often wish to refer to certain
pairs of the angles formed. In order to be precise, we need a definition.

\begin{definition}
A line $l$ (Fig.~8--9) that intersects lines $l_1$ and $l_2$ in distinct
points is called a \emph{transversal} of $l_1$ and $l_2$. We also say that
$l_1$ and $l_2$ are cut by the transversal $l$. Angles 3, 4, 5, and 6 are
called \emph{interior angles}. Angles 1, 2, 7, and 8 are called \emph{exterior
angles}. These eight angles are grouped into four pairs referred to as pairs
of \emph{corresponding angles}. Angles 1 and 5 form such a pair of
corresponding angles, as do also 2 and 6, 3 and 7, and 4 and 8. Angles such as
3 and 5 are called \emph{alternate angles}. Other pairs of alternate angles
are 2 and 8, 1 and 7, and 4 and 6. We speak of the pair 3 and 5 as
\emph{alternate interior angles}. We call the pair 2 and 8 \emph{alternate
exterior angles}. We observe that \emph{alternate} angles have the points of
their interiors on \emph{opposite} sides of $l$, while \emph{corresponding}
angles have their interiors on the \emph{same} side of $l$. In a pair of
corresponding angles, one angle is always exterior and one interior.
\end{definition}

\FIGVIIININETEN

\begin{theorem}
Let $l_1$ and $l_2$ be lines cut by line $l$. If $l_1 \pll l_2$, then two
alternate interior angles are congruent to each other, and conversely, if two
alternate interior angles are congruent, then $l_1 \pll l_2$
\textup{(Fig.~8--10)}.
\end{theorem}

\emph{Proof.} The last part of the theorem follows at once from the first
proof of Theorem 8--1, for if $\ang{CAB} \cong \ang{ABD}$, then
$\ang{ABD} \cong \ang{EAF}$, and $l_1 \pll l_2$.

The first part of the theorem requires the parallel postulate, and may be
proved as follows. Choose $G$, located in the interior of $\ang{DBE}$, such
that $\ang{EAG} \cong \ang{ABD}$. Then line $AG \pll l_2$. (Why?) But $A$ is
also on $l_1$, and $l_1 \pll l_2$. Hence the line through $A$ and $G$ is the
same as $l_1$, so $G$ is on $l_1$. Now,
$\ang{CAB} \cong \ang{EAG} \cong \ang{ABD}$. Thus the alternate interior
angles are congruent.

As important consequences of Theorem 8--4, we state the following corollaries,
which you can easily prove.

\begin{corollary}
Let $l$ cut $l_1$ and $l_2$ at distinct points. If $l_1 \pll l_2$, then
alternate exterior angles are congruent, and conversely.
\end{corollary}

\begin{corollary}
Let $l$ cut $l_1$ and $l_2$ at distinct points. If $l_1 \pll l_2$, then
corresponding angles are congruent, and conversely.
\end{corollary}

\exercisegroup{8--6}
\begin{exercisecols}
\begin{exlist}
\item Prove that if a line is perpendicular to one of two parallel lines, it
  is perpendicular to the other.
\item Prove that if $l \pll m$, $n \perp l$, $p \perp m$, and $n \neq p$, then
  $n \pll p$.
\item In Fig.~8--11, $\sg{AO} = \sg{CO}$, and $\sg{BO} = \sg{OD}$. Prove that
  $AB \pll CD$.
\exfig{\FIGVIIIELEVEN}
\item Show that $\sg{BC} = \sg{AD}$, and $BC \pll AD$ in Fig.~8--11.
\item Show that if $AB \cong CD$, and $AB \pll CD$, then $BO \cong OD$, and
  $BC \pll AD$ in Fig.~8--11.
\item If $l \pll m$, and $A$ and $B$ are any two points on $l$, show that the
  perpendicular distances from $A$ and $B$ to $m$ are equal. We may state this
  result in the form: if $l \pll m$, then $l$ and $m$ are ``everywhere
  equidistant.'' State and prove the converse of this theorem. State also the
  contrapositive.
\end{exlist}
\end{exercisecols}

%========================================================= 8-6
\section*{8--6\quad ANGLE SUMS}
\addcontentsline{toc}{section}{8--6\quad Angle sums}

One of the most interesting consequences of the parallel postulate is a
theorem relating to the angles of a triangle. We need a definition.

\begin{definition}
If $\ang{AOB}$ and $\ang{BOC}$ are adjacent angles (Fig.~8--12), and either
$\ang{AOC}$ is a straight angle, or $OB$ is in the interior of $\ang{AOC}$,
then we say that $\ang{AOC}$ is a \emph{sum} of $\ang{AOB}$ and $\ang{BOC}$.
Likewise, if $\ang{A'O'B'} \cong \ang{AOB}$, and
$\ang{B''O''C''} \cong \ang{BOC}$, then we say that $\ang{AOC}$ is a
\emph{sum} of $\ang{A'O'B'}$ and $\ang{B''O''C''}$.
\end{definition}

\FIGVIIITWELVE

\begin{remark}
At this stage in our geometry, we cannot yet ``add'' angles by adding their
measures in degrees, because we have not as yet described how angles are
measured. The measurement of angles is more complicated than measurement of
length, and will be described in a later chapter.
\end{remark}

Draw a figure to illustrate the last part of Definition 8--4.

\begin{definition}
Two angles are said to be \emph{complementary} if their sum is a right angle.
\end{definition}

\begin{definition}
Two angles are said to be \emph{supplementary} if their sum is a straight
angle.
\end{definition}

\exercisegroup{8--7}
\begin{exercisecols}
\begin{exlist}
\item Draw an obtuse angle. Draw an angle that is supplementary to this angle
  and is not adjacent to it. Draw one that is adjacent and supplementary.
\item Draw an acute angle. Draw an angle that is complementary to this angle
  and is adjacent to it. Draw one that is not adjacent, but is complementary,
  to the acute angle.
\item Prove that if two acute angles are congruent, then their complements are
  congruent.
\item Prove that if the sides of one angle are parallel, respectively, to the
  sides of a second angle, then either the angles are congruent or they are
  supplementary.
\end{exlist}
\end{exercisecols}

\begin{theorem}
The sum of the angles of a triangle is a straight angle.
\end{theorem}

\FIGVIIITHIRTEEN

\emph{Proof.} In Fig.~8--13, let $l$ be the line through $A$ and parallel to
the side $BC$ of $\tri{ABC}$. Let $D$ be a point on $l$ that is on the
opposite side of $AB$ from $C$. Let $E$ be a point of $l$ on the same side of
$AB$ as $C$. Then $\ang{DAB} \cong \ang{ABC}$, and
$\ang{EAC} \cong \ang{ACB}$. The straight angle $DAE$ is the sum of angles
$A$, $B$, and $C$.

\exercisegroup{8--8}
\begin{exercisecols}
\begin{exlist}
\item Show that the acute angles of a right triangle are complementary.
\item Two angles of a triangle are $40^\circ\ 17'\ 30''$ and
  $78^\circ\ 31'\ 38''$, respectively. Determine the number of degrees in the
  third angle.
\item How many degrees are there in each angle of an equilateral triangle?
\item Show that a median of an equilateral triangle divides the triangle into
  two right triangles, each of which has acute angles of $30^\circ$ and
  $60^\circ$.
\item Prove that in a right triangle, with acute angles of $30^\circ$ and
  $60^\circ$, the side opposite the $30^\circ$ angle is half as long as the
  hypotenuse.
\item In $\tri{ABC}$, an exterior angle at $B$ is $124^\circ$. $\ang A$ is
  $58^\circ$. How many degrees in $\ang C$?
\item Prove that if two angles of one triangle are congruent, respectively, to
  two angles of another triangle, then the third angles are congruent.
\item In Fig.~8--14, $AB \cong A'B'$, $\ang A \cong \ang{A'}$, and
  $\ang C \cong \ang{C'}$. Prove that $\tri{ABC} \cong \tri{A'B'C'}$.
\exfig{\FIGVIIIFOURTEEN}
\item Is the following true? Two triangles are congruent to each other if one
  side and two angles of the one triangle are congruent to one side and two
  angles of the other.
\item Prove that two right triangles are congruent if the hypotenuse and an
  acute angle of one triangle are congruent, respectively, to the hypotenuse
  and acute angle of the other.
\item In $\tri{ABC}$ (Fig.~8--15), $\ang A$ is right, and $AD \perp BC$. Prove
  that $\ang{BAD} \cong \ang{ACD}$, and $\ang{ABD} \cong \ang{DAC}$.
\exfig{\FIGVIIIFIFTEEN}
\item[\stex 12.] In Fig.~8--16, $D$ is the midpoint of $BC$, and
  $BD \cong AD$. Prove that $\ang{BAC}$ is right.
\exfig{\FIGVIIISIXTEEN}
\item[\stex 13.] Prove that two right triangles are congruent if the
  hypotenuse and one side of one triangle are congruent to the hypotenuse and
  one side of the other.
\end{exlist}
\end{exercisecols}

%========================================================= 8-7
\section*{8--7\quad POLYGONS}
\addcontentsline{toc}{section}{8--7\quad Polygons}

In this and the remaining sections, we will be studying certain polygons.
First, we need to recall the definition of a polygonal path.

A \emph{polygonal path} (Fig.~8--17) is determined by a number of points,
$P_1, P_2, \ldots, P_n$, called the \emph{vertices}, given in a definite
order. The path is the set of points on the segments
$P_1P_2, P_2P_3, \ldots, P_{n-1}P_n$. The point $P_1$ is called the
\emph{beginning} of the path, and $P_n$ is called the \emph{end}. The segments
$P_1P_2$, etc., are called the \emph{sides} of the path.

\FIGVIIISEVENTEENEIGHTEEN

\begin{remark}
There is no requirement that the sides of a polygonal path not intersect.
Hence the path can be quite complicated. See Fig.~8--18 for a six-sided path.
\end{remark}

\exercisegroup{8--9}
\begin{exercisecols}
\begin{exlist}
\item Draw a polygonal path with five sides that does not cross itself. Draw
  one that does.
\item If the length of a polygonal path is the sum of the lengths of its
  sides, how does the length of a polygonal path $P_1, \ldots, P_n$ compare
  with $\sg{P_1P_n}$?
\end{exlist}
\end{exercisecols}

\begin{definition}
A \emph{polygon} is a polygonal path whose beginning and end coincide. If the
vertices are all different and no two sides have a point in common (other than
a vertex of two adjacent sides), then the polygon is called \emph{simple}
(Fig.~8--19).
\end{definition}

\FIGVIIININETEEN

Each simple polygon separates the points of the plane, other than the polygon
itself, into two sets called the \emph{interior} and \emph{exterior} of the
polygon. If $A$ is a point of the interior, and $B$ is a point of the
exterior, every polygonal path connecting $A$ to $B$ contains at least one
point of the polygon. If $A$ and $B$ both belong to either the interior or
exterior, then there exist polygonal paths joining $A$ to $B$ that contain no
point of the polygon. The proofs of these statements are quite difficult, and
will not be given here. It would be instructive to carry out the proofs in the
case in which the polygon is a triangle. The student interested in doing so
should restudy Section 3--6.

\emph{*Problem.} The theorem that a simple polygon separates the plane is
powerful and has surprising applications. It can be used to show that the
well-known puzzle below has no solution.

It is desired to connect each of three houses $A$, $B$, and $C$ with three
outlets, one for gas $G$, one for water $W$, and one for electricity $E$.

\VIIIutilities

Is it possible to run the lines, one from each house to each outlet, in such a
way that no lines cross? In other words, can polygonal paths from $A$ to $G$,
$A$ to $W$, etc., be drawn so that no two paths cross?

\begin{definition}
Vertices of a polygon that are endpoints of the same side are called
\emph{adjacent} vertices (Fig.~8--20). Angles whose vertices are adjacent to
each other are called \emph{consecutive} angles. A \emph{diagonal} of a
polygon is a segment joining any two nonadjacent vertices. Sides with a common
vertex are called \emph{adjacent} sides.
\end{definition}

\FIGVIIITWENTY

\begin{definition}
A simple polygon is called \emph{convex} if, whenever $A$ and $B$ are points
of the polygon, every point of the segment $AB$ lies either on the polygon or
in its interior (see Fig.~8--21).
\end{definition}

\FIGVIIITWENTYONE

\begin{remark}
Unless otherwise stated, we will confine our attention to simple convex
polygons.
\end{remark}

\begin{definition}
If all sides of a polygon have the same length, the polygon is
\emph{equilateral}. If all angles are congruent, it is \emph{equiangular}. A
convex polygon that is equilateral and equiangular is called \emph{regular}.
Polygons are also classified according to the number of sides. Polygons of 3,
4, 5, 6, 7, 8, 9, and 10 sides are called \emph{triangles},
\emph{quadrilaterals}, \emph{pentagons}, \emph{hexagons}, \emph{heptagons},
\emph{octagons}, \emph{nonagons}, and \emph{decagons}, respectively. A polygon
of $n$ sides will be called an \emph{$n$-gon}.
\end{definition}

\exercisegroup{8--10}
\begin{exercisecols}
\begin{exlist}
\item What do we call a regular polygon of three sides? of four sides?
\item Draw a convex polygon of five sides. Draw a quadrilateral that is not
  convex. Draw a polygon that is not simple.
\item[\stex 3.] Is it true that every interior point of a convex polygon is
  also interior to every angle of the polygon?
\item[4.] We can divide a convex polygon into triangles by drawing the
  diagonals from any vertex. Into how many triangles does this procedure
  divide a convex polygon of four sides? of ten sides? of $n$ sides?
\item[\stex 5.] How many diagonals has a convex polygon of four sides? five
  sides? ten sides? $n$ sides?
\item[6.] Show that an equiangular polygon need not be equilateral.
\item[\stex 7.] Find an example of an equilateral, equiangular polygon that is
  not regular.
\end{exlist}
\end{exercisecols}

A diagonal of a convex quadrilateral divides the quadrilateral into two
triangles. Since the sum of the angles in each triangle is a straight angle,
we shall say that \emph{the sum of the (four) angles of the quadrilateral is
two straight angles}. The two diagonals from one vertex of a convex pentagon
divide the pentagon into three triangles, and we say that \emph{the sum of the
(five) angles of the pentagon is three straight angles}.

When we have defined degree measure for angles, we will see that the sum of
the degree measures of the angles of a convex quadrilateral is
$2 \cdot 180^\circ = 360^\circ$, of a pentagon,
$3 \cdot 180^\circ = 540^\circ$, and in general for a polygon of $n$ sides,
the sum of degree measures is $(n - 2) \cdot 180^\circ$.

Do not misunderstand the statement that the sum of the angles of a pentagon is
three straight angles. This is using the word ``sum'' in a different sense
than when we speak of the sum of \emph{two angles}. Recall that when we add
two angles, the sum is less than or congruent to a straight angle.

\FIGVIIITWENTYTWOTHREE

We observe that the sum of the five interior angles and the five indicated
exterior angles in Fig.~8--22 is precisely five straight angles. Since the
interior angles account for three straight angles, the sum of the exterior
angles must be two straight angles. Similar reasoning shows that for any
convex polygon of $n$ sides, the sum of the $n$ exterior angles (one at each
vertex) is exactly two straight angles (Fig.~8--23).

\exercisegroup{8--1}
\begin{exercisecols}
\begin{exlist}
\item A quadrilateral has three $95^\circ$ angles. How many degrees are there
  in the fourth angle?
\item How many degrees are there in each angle of a regular hexagon?
\item How many degrees are there in each exterior angle of a regular pentagon?
\item A convex quadrilateral has three angles of the same size. The fourth
  angle is a $120^\circ$ angle. How many degrees are there in each of the
  other angles?
\item Explain why it is impossible for a convex polygon to have as many as six
  angles smaller than $120^\circ$.
\item[\stex 6.] What is the sum of the angles at the vertices of a regular
  five-pointed star? (The points of the star are vertices of a regular
  pentagon.)
\exfig{\FIGVIIITWENTYFOUR}
\item[7.] An exterior angle of a regular polygon is $20^\circ$. How many sides
  has the polygon?
\item[8.] Pentagon $ABCDE$ is regular. Prove that a line $l$, bisecting the
  angle at $A$, is perpendicular to side $CD$.
\item[9.] How many sides does a regular polygon have if each angle is
  $120^\circ$? $140^\circ$? $160^\circ$? $179^\circ$?
\item[10.] Make up some theorems of your own relating to regular pentagons or
  regular hexagons.
\end{exlist}
\end{exercisecols}

%========================================================= 8-8
\section*{8--8\quad PARALLELOGRAMS}
\addcontentsline{toc}{section}{8--8\quad Parallelograms}

The following definitions make certain familiar terms precise.

\begin{definition}
A \emph{parallelogram} is a quadrilateral having two pairs of parallel sides.
A \emph{rectangle} is a parallelogram with a right angle. A \emph{square} is a
rectangle with two congruent perpendicular sides. A \emph{rhombus} is an
equilateral parallelogram.
\end{definition}

\FIGVIIITWENTYFIVE

Our knowledge of parallels and these definitions make it possible to prove
many results, as the following exercises will demonstrate.

\exercisegroup{8--12}
\noindent The exercises below present many interesting theorems concerning
polygons. Give proofs for these theorems.

\begin{exercisecols}
\begin{exlist}
\item A diagonal of a parallelogram divides it into two congruent triangles.
\item Opposite sides of a parallelogram are congruent to each other.
\item The diagonals of a parallelogram bisect each other.
\item Opposite angles of a parallelogram are congruent.
\item Consecutive angles of a parallelogram are supplementary to each other.
\item Every angle of a rectangle is right.
\item All sides of a square are congruent.
\item The diagonals of a rectangle are congruent.
\item The diagonals of a square are perpendicular to each other.
\item The diagonals of a rhombus are perpendicular.
\item If the diagonals of a parallelogram are congruent, the parallelogram is
  a rectangle.
\item If the diagonals of a rectangle are perpendicular, the rectangle is a
  square.
\item The quadrilateral formed by connecting midpoints of the sides of a
  parallelogram is also a parallelogram.
\item The parallelogram formed by connecting midpoints of the sides of a
  square is also a square.
\item If the diagonals of a quadrilateral bisect each other, the quadrilateral
  is a parallelogram.
\item If a quadrilateral has a pair of sides congruent and parallel to each
  other, it is a parallelogram.
\item A parallelogram has one angle of $70^\circ$. Give the number of degrees
  in each angle.
\item One angle of a parallelogram is twice as large as a second angle.
  Determine the number of degrees in each angle of the parallelogram.
\item A diagonal of a rhombus bisects two angles of the rhombus.
\item If a rectangle is not a square, a diagonal does not bisect an angle of
  the rectangle.
\item If each angle of a quadrilateral is congruent to the angle opposite it,
  then the quadrilateral is a parallelogram.
\item Is a rhombus a regular polygon?
\item For every quadrilateral, it is true that either the bisectors of two
  consecutive angles of the quadrilateral are perpendicular, or the
  quadrilateral is not a parallelogram.
\item If a diagonal of a parallelogram bisects one angle of the parallelogram,
  it bisects a second angle, and the parallelogram is a rhombus.
\item If a parallelogram is not a rhombus, then the four bisectors of the
  angles of the parallelogram are not concurrent.
\item If a parallelogram is not a rhombus, then the four bisectors of the
  angles intersect to form a rectangle.
\item The four bisectors of the angles of a rectangle that is not a square
  intersect to form a square.
\item How many degrees are there in the angle formed by the two diagonals
  drawn from one vertex of a regular pentagon?
\item In Fig.~8--26, $DE \pll AC$, and $DF \pll AB$. Use the figure to prove
  that the sum of the angles of $\tri{ABC}$ is a straight angle.
\exfig{\FIGVIIITWENTYSIX}
\item In the regular hexagon $ABCDEF$, prove that diagonal $FD$ is
  perpendicular to side $DC$ (Fig.~8--27).
\exfig{\FIGVIIITWENTYSEVEN}
\item Prove that $\sg{FC} = 2 \cdot \sg{DC}$ in Fig.~8--27. Show that the
  diagonals $FC$ and $AD$ bisect each other.
\end{exlist}
\end{exercisecols}

%========================================================= Review of Chapter 8
\section*{REVIEW OF CHAPTER 8}
\addcontentsline{toc}{section}{Review of Chapter 8}

This chapter has dealt with consequences of the parallel postulate, by far the
most famous of all the postulates of Euclid. (Euclid did not state it exactly
as we have given it.) It was through the careful investigation of this
postulate that non-Euclidean geometries were invented early in the 19th
century. This investigation also led indirectly to the first set of correct
postulates from which the theorems of Euclidean geometry could be rigorously
proved. These postulates, presented about 1900 by the German mathematician
David Hilbert, are equivalent to the ones that we have listed in this text. We
list below some of the most important consequences of the parallel postulate
shown in this chapter.

\begin{quote}
If two lines are parallel to one line, they are parallel to each other.

When parallel lines are cut by a line, alternate interior angles are
congruent.

The sum of the angles of any triangle is a straight angle.

The sum of the interior angles of any convex polygon is $(n - 2)$ straight
angles.
\end{quote}

Several new terms were defined in this chapter. In addition, precise
definitions of some familiar concepts were formulated. (For example, though we
have referred freely to squares and rectangles in earlier chapters, they had
not yet been exactly defined. Indeed, it is not even possible to prove that
squares and rectangles exist until the parallel postulate is available.) You
should know the definitions given in this chapter for the following:

\begin{center}
\begin{tabular}{@{}l@{\hspace{3.2em}}l@{}}
alternate interior angles & octagon \\
corresponding angles      & $n$-gon \\
polygon                   & regular polygon \\
simple polygon            & parallelogram \\
convex polygon            & rectangle \\
parallel lines            & square \\
pentagon                  & rhombus \\
hexagon                   & diagonal \\
equilateral polygon       & equiangular polygon \\
\end{tabular}
\end{center}

%--------------------------------------------------------- Algebra Review
\section*{ALGEBRA REVIEW}
\addcontentsline{toc}{section}{Algebra Review (Chapter 8)}

\begin{exlist}
\item Two angles of a triangle are congruent, and the measure of each exceeds
  that of the third angle by $10^\circ$. Determine the number of degrees in
  each angle of the triangle.
\item In $\tri{ABC}$, the sum of $\ang B$ and an exterior angle at $A$ has a
  measure of $170^\circ$. The sum of this exterior angle and $\ang C$ is
  $160^\circ$. Determine the measure of each angle of the triangle.
\item An exterior angle of a triangle has a measure of $117^\circ$. One angle
  of this triangle has a measure of $42^\circ$. Determine the degree measure
  of each angle of the triangle.
\item Prove that there is no $\tri{ABC}$ such that the measure of an exterior
  angle at $A$ is $103^\circ$ and an exterior angle at $B$ is $74^\circ$.
\item Diagonal $AC$ of parallelogram $ABCD$ forms with sides $AB$ and $AD$ two
  angles whose measures differ by $10^\circ$. The measure of $\ang B$ is
  $110^\circ$. Determine the measure of $\ang A$.
\item Two angles of a triangle are complementary. Prove that the measure of
  the third angle is $90^\circ$.
\item The measure of the vertex angle of an isosceles triangle is $84^\circ$.
  Find the measure of each angle of the triangle.
\item The first angle of a triangle has a measure three-fourths of the second.
  The first angle is also one-third of the sum of the other two. Compute each
  angle.
\item In Fig.~8--28 the letters represent the degree measures of the indicated
  angles. Prove that $x + y + t = r + s + z$.
\exfig{\FIGVIIITWENTYEIGHT}
\item In Fig.~8--29, $AB \cong BC$, $BE \cong BD$, and the measure of
  $\ang{CBD}$ is $40^\circ$. Determine the measure of $\ang{ADE}$.
\exfig{\FIGVIIITWENTYNINE}
\end{exlist}

%========================================================= Review, Chs 1-8
\section*{REVIEW EXERCISES,\\CHAPTERS 1 THROUGH 8}
\addcontentsline{toc}{section}{Review Exercises, Chapters 1 through 8}

\begin{exlist}
\item Give truth values of $\al$, $\be$, and $\ga$ for which a statement of the
  form ($\al$ or $\be$) and $\ga$ is true; is false.
\item Can you find truth values of $\al$, $\be$, and $\ga$ for which ($\al$ or
  $\be$) and $\ga$ is false, and $\al$ or ($\be$ and $\ga$) is true? Can you
  find truth values for which ($\al$ or $\be$) and $\ga$ is true, and $\al$ or
  ($\be$ and $\ga$) is false? Is ($\al$ or $\be$) and
  $\ga \imp \al$ or ($\be$ and $\ga$) true? Is $\al$ or ($\be$ and $\ga$)
  $\imp$ ($\al$ or $\be$) and $\ga$ true?
\item Formulate an ``if\ldots, then\ldots'' statement equivalent to (a) no
  fishermen lie, (b) only fishermen lie, (c) fishermen necessarily lie.
\item Give the contrapositive of each of your implications in Exercise 3.
\item One set of points contains, among others, the points $A$, $B$, $C$, and
  $D$. A second set of points contains $A$, $B$, and $D$, but does not contain
  the point $C$. Prove that these two sets of points cannot both be lines.
\item Point $C$ is in the interior of $\ang{AOB}$. Is every point of ray $OC$
  in the interior of this angle?
\item Define \emph{interior of a triangle} in at least two ways.
\item Explain what is meant by saying that \emph{congruence of segments} has
  the \emph{reflexive} property; the \emph{symmetric} property; the
  \emph{transitive} property.
\item Explain why one of these statements is true, and the other false:\\
  (a) $AB \cong CD \imp AB = CD$, (b) $AB = CD \imp AB \cong CD$.
\end{exlist}

\begin{exlist}
\item[10.] Figure 8--30 indicates that $AB \cong A'B'$, and $ABC$. Using only
  the definitions of \emph{greater than} and \emph{less than}, does the figure
  indicate that $AC > A'B'$, or that $A'B' < AC$?
\exfig{\FIGVIIITHIRTY}
\item[11.] The length of $AB$, measured with $CD$ as unit, is 17.23. What is
  Archimedes' integer?
\item[12.] The length of $CD$ in Exercise 11, as measured by $EF$, is 10. If
  $AB$ is measured by $EF$, what is its length? What is Archimedes' integer in
  this case?
\item[13.] Express the rational number $\tfrac{3}{19}$ as a repeating decimal.
  Find the rational number whose decimal representation is the repeating
  decimal $0.272727\ldots$; $0.027027027\ldots$.
\item[14.] State the S.A.S. postulate for triangles.
\item[15.] Two lines that intersect to form congruent adjacent angles are said
  to be perpendicular to each other. Does this definition in itself prove that
  there is a pair of perpendicular lines in the plane? How do we know that
  right angles exist?
\item[16.] Are there points $A$, $B$, and $C$, such that $\sg{AB} = 5$,
  $\sg{BC} = 7$, and $\sg{AC} = 13$?
\item[17.] What can be said about the position of the points $R$, $S$, and
  $T$, if $\sg{RS} = 8$, $\sg{ST} = 19$, and $\sg{RT} = 11$?
\item[18.] Consider a circle whose center is $A$ and which contains the point
  $B$.
  \begin{enumerate}[label=(\alph*), leftmargin=1.8em, topsep=2pt, itemsep=1pt]
    \item If $\sg{AX} > \sg{AB}$, what can be said about point $X$?
    \item If $P$ and $Q$ are on the circle, how do we refer to $\ang{PAQ}$? to
      segment $PQ$?
    \item If $O$ and $C$ are on the circle, how do we refer to $\ang{AOC}$?
    \item If $R$ and $S$ are points of the circle, and $A$ is on the segment
      $RS$, what do we call segment $RS$?
  \end{enumerate}
\item[19.] What does it mean to say that two rays have the same direction?
\item[\stex 20.] Without using the parallel postulate or any of its
  consequences, prove the following theorem.
  \begin{quote}
  \emph{If in triangles $ABC$ and $A'B'C'$, $AB \cong A'B'$,
  $\ang B \cong \ang{B'}$, and $\ang C \cong \ang{C'}$, then
  $\tri{ABC} \cong \tri{A'B'C'}$.}
  \end{quote}
\item[\stex 21.] Without using the parallel postulate, prove the following
  theorem, which relates to Fig.~8--31.
  \begin{quote}
  \emph{If $\ang{CAB}$ and $\ang{ABD}$ are right, $CA \cong DB$, and $M$ and
  $N$ are midpoints of $CD$ and $AB$, respectively, then $MN$ is perpendicular
  both to $AB$ and to $CD$.}
  \end{quote}
\exfig{\FIGVIIITHIRTYONE}
\end{exlist}

\begin{exlist}
\item[\stex 22.] Carry through the proofs of Exercises 20 and 21 using the
  parallel postulate.
\item[\stex 23.] Euclid's formulation of the parallel postulate is different
  from ours. The following statement is essentially the postulate that Euclid
  used:
  \begin{quote}
  \emph{If the points $X$ and $Y$ are on one side of line $AB$, and if the sum
  of $\ang{XAB}$ and $\ang{YBA}$ is less than a straight angle, then the rays
  $AX$ and $BY$ intersect.}
  \end{quote}
  Figure 8--32 illustrates Euclid's parallel postulate. (a) Prove Euclid's
  parallel postulate as a theorem, using our parallel postulate. (b) Assuming
  the truth of Euclid's parallel postulate, prove ours as a theorem.
\exfig{\FIGVIIITHIRTYTWO}
\end{exlist}

\end{document}
